% HD-F1-T40. GPT-6 Astra Pro, /root/transport; 22 November 2026 UTC.
\section[HD-F1-T40: a split orthogonal rank cutoff]{HD-F1-T40: an effective split orthogonal rank bound}
\label{sec:f1-split-orthogonal-rank-cutoff}

With \(L_0=3\cdot10^{30}-3\), every original remaining
T2-OTHER source whose almost simple projective socle is
\(P\Omega^+_{2r}(q)\) satisfies
\[
 \boxed{\qquad r\le4L_0^2+8<10^{62}.\qquad}
 \tag{F1T40.1}
\]
The retained field bound is \(q\le27\); both characteristics
are included. The proof uses the actual bounded-level
permutation module, with a commuting symplectic action
constructed by finite Fourier transforms. It requires no
additional duality theorem. Nonsplit even-dimensional and
odd-dimensional orthogonal sources are not asserted to satisfy
(F1T40.1).

\paragraph{1. The actual natural quotient and its automorphisms.}
Retain the original \(G\subset GL_K(V)\),
\(m=\dim_KV\ge27\), scalar centre \(Z\), perfect layer
\(L\), and the actual T35--T36 surjection
\[
 \beta:H=\Omega^+(A)\longrightarrow L,
 \qquad A=X\oplus X^*,\quad\dim X=r,
 \quad Q_A(x,y)=x^ty.
 \tag{F1T40.2}
\]
Suppose \(r\ge4L_0^2+9\). All covering exceptions and
triality ranks lie below this range. T36 gives the actual
lifted automorphism square for every original \(g\in G\).
Malle--Testerman Theorem 24.24, with its full split,
simply-laced proof, identifies the projective automorphism
as inner, diagonal, field and graph. Each factor has an
actual semilinear-similitude representative on this same \(A\).
For clarity, the diagonal factor can be written explicitly.
For simple roots
\(\varepsilon_i-\varepsilon_{i+1}\) and
\(\varepsilon_{r-1}+\varepsilon_r\), prescribe nonzero
root scalars \(c_1,\ldots,c_r\). Set
\[
 a_r=1,\qquad a_i=\prod_{j=i}^{r-1}c_j,
 \qquad\lambda=c_{r-1}/c_r,
 \qquad d=\operatorname{diag}(a_i,\lambda a_i^{-1}).
 \tag{F1T40.3}
\]
This natural similitude has the prescribed simple-root
multipliers. The nontrivial diagram action is represented
by interchanging the last hyperbolic pair; it exchanges
the two terminal simple roots. Field actions are entrywise
Frobenius, and inner actions are natural matrices. Thus
\[
 \alpha_g(h)=a h^{p^e}a^{-1},\qquad
 a\in\operatorname{GO}(A),\qquad
 \beta\alpha_g=c_g\beta.
 \tag{F1T40.4}
\]
The equality with the actual lifted action is exact: two
lifts with the same projective action differ by a
homomorphism from the perfect \(H\) to its centre, and
that homomorphism is trivial. The natural central kernel
was already retained in T35--T36.

The group \(H\) is generated by its root subgroups.
In the displayed hyperbolic basis their nonidentity
matrices have the forms
\[
 I+tE_{ij}-tE_{r+j,r+i},\qquad
 I+tE_{i,r+j}-tE_{j,r+i},
 \tag{F1T40.5}
\]
and the corresponding negative-root matrices, with
\(i\ne j\) and \(t\ne0\). Each has moving rank two.
For every \(g\notin Z\), at least one such actual
root element \(u\) satisfies
\(\beta(\alpha_g(u)u^{-1})\notin Z(L)\).
Otherwise the original projective action would fix all
root generators of the simple socle and centralize it,
contrary to its trivial centralizer in the actual almost
simple projective group. Put
\[
 h=\alpha_g(u)u^{-1},\qquad
 \beta(h)=[g,\beta(u)].
 \tag{F1T40.6}
\]
Equation (F1T40.4) preserves natural rank, so
\(\operatorname{rank}_A(h-I)\le4\).
The rank is nonzero and even. In odd characteristic this
is the fixed-space parity in the orthogonal part of
Guralnick--Larsen--Tiep, Lemma 3.3. In even characteristic
it is the quasideterminant kernel identity of
Guralnick--Tiep, Lemma 5.8(ii), also used in that lemma.
These named published parity inputs give
\[
                  2\le\operatorname{rank}_A(h-I)\le4.
 \tag{F1T40.7}
\]
No original inertia has been replaced or rescaled.

\paragraph{2. Actual permutation levels and their Fourier coordinates.}
The same published Lemma 3.3 gives
\(\zeta_A=\tau_A\) on \(H\), where \(\tau_A\) is
its permutation character on \(A\). We use this equality
through an actual module isomorphism, as follows.
Decompose the two genuine modules by their averaged
central idempotents. For each irreducible \(W\), the
spaces \(\operatorname{Hom}_H(W,\zeta_A)\) and
\(\operatorname{Hom}_H(W,\tau_A)\) have equal dimensions,
by the character inner product. Choose a linear isomorphism
between them and compose it with the actual evaluation
isomorphisms \(W\otimes\operatorname{Hom}_H(W,-)\).
Their direct sum gives an invertible \(H\)-intertwiner
\(J:\zeta_A\longrightarrow\mathbf C[A]\).
All these Hom spaces are the images of the finite averaging
maps on ordinary linear Hom spaces. Consequently
\[
 \tau_A\oplus\zeta_A\simeq\mathbf C[A]\otimes\mathbf C^2
 \tag{F1T40.8}
\]
by the map which is the identity on the first summand and
\(J\) on the second. A constituent's minimal T36 level
is therefore precisely its minimal permutation tensor level.
For each actual constituent select a nonzero coordinate
projection of its tensor inclusion under (F1T40.8).
Irreducibility makes that projection an injection into
\[
          \mathcal R_k=\mathbf C[A^k],\qquad
                    1\le k\le L_0.
 \tag{F1T40.9}
\]
This is a comparison of the given coefficient constituent,
not a substitution by an ambient tensor factor.

Fix a nontrivial additive character \(\psi\) of
\(\mathbf F_q\). Write the \(k\) columns of \(A^k\)
as pairs of \(r\)-by-\(k\) matrices \((x,y)\).
The actual partial Fourier basis is
\[
 e_{x,z}=q^{-rk/2}\sum_y\psi\bigl(\operatorname{Tr}(y^tz)\bigr)
                                     \delta_{(x,y)}.
 \tag{F1T40.10}
\]
For each row set \(F_i=(x_i,z_i)\) in the symplectic
space \(U=\mathbf F_q^k\oplus\mathbf F_q^k\), with
\[
 b_U((x,z),(x',z'))=z\cdot x'-x\cdot z'.
 \tag{F1T40.11}
\]
Thus (F1T40.10) identifies \(\mathcal R_k\) with
\(\mathbf C[\operatorname{Hom}(X,U)]\), of dimension
\(D=q^{2rk}\). A Levi element \(a\in SL(X)\)
sends \(e_{x,z}\) to \(e_{ax,az}\), the corresponding
change of \(X\)-coordinates in both sets of rows. For an alternating
matrix \(B\), where \(B^t=-B\) and \(B_{ii}=0\), the
actual lower Siegel element \((x,y)\mapsto(x,y+Bx)\) acts by
\[
 e_F\longmapsto
 \psi\!\left(-\sum_{i<j}B_{ij}b_U(F_i,F_j)\right)e_F.
 \tag{F1T40.12}
\]
The minus sign follows from replacing \(y\) by \(y-Bx\)
in (F1T40.10); it is retained in odd characteristic too.
Interchanging the \(i\)-th hyperbolic pair acts on its
row by
\[
 \mathcal F_U e_x=q^{-k}\sum_y\psi(b_U(x,y))e_y,
                  \qquad \mathcal F_U^2=I.
 \tag{F1T40.13}
\]

Only paired such interchanges are needed inside \(H\).
They do lie in the actual root-generated group. For a pair
\(i,j\), let \(n_{ij}\) and \(n_{i+j}\) be the root
\(SL_2\) Weyl matrices
\(u_\alpha(1)u_{-\alpha}(-1)u_\alpha(1)\)
for \(\varepsilon_i-\varepsilon_j\) and
\(\varepsilon_i+\varepsilon_j\). Matrix multiplication
in their four-dimensional hyperbolic block gives
\[
 n_{ij}n_{i+j}=\ell_{ij}w_iw_j,
 \tag{F1T40.14}
\]
where \(w_i,w_j\) interchange the respective pairs and
\(\ell_{ij}\) is the Levi element with \(-1\) in the
two indicated coordinates and \(+1\) elsewhere.
Its determinant on \(X\) is one, so \(\ell_{ij}\in H\).
This proves (F1T40.14) in both characteristics, without
using a disconnected orthogonal operation as an \(H\)-operator.

The full \(S=\operatorname{Sp}(U)\) acts by the actual
permutations \(e_F\mapsto e_{sF}\). Equations
(F1T40.11)--(F1T40.13) show that these commute with the
Levi, Siegel and paired Fourier operators. These include
all root-group generators: the difference roots lie in
\(SL(X)\), one sign of sum roots lies in the Siegel group,
and paired interchanges conjugate them to the other sign.
Thus they give a genuine commuting \(H\times S\)-action.
No scalar extension or additional representation theorem
is needed for this action.

\paragraph{3. Actual isotropic codes and the paired-Fourier proof.}
For an isotropic subspace \(N\subset U\) of dimension
\(j\), put \(U_N=N^\perp/N\) and use the normalized map
\[
 i_N e_{\bar F}=q^{-rj/2}\sum_{\Delta:X\to N}e_{F+\Delta}.
 \tag{F1T40.15}
\]
It is an isometry whose image \(C_N\) consists of the
functions supported on \(\operatorname{Hom}(X,N^\perp)\)
and constant on its \(\operatorname{Hom}(X,N)\)-cosets.
The phases (F1T40.12) descend to the quotient, the Levi
operators commute with the sums, and finite Fourier
orthogonality checks (F1T40.13) on each row, with its
normalization. Hence (F1T40.15) is an actual \(H\)-map.
A symplectic basis of \(U_N\) identifies it with the
standard space of dimension \(2(k-j)\); the corresponding
basis permutation intertwines every displayed generator.
Inverting the actual partial Fourier map (F1T40.10) then gives
\[
                  C_N\simeq\mathbf C[A^{k-j}],
                  \qquad\dim C_N=q^{2r(k-j)}.
 \tag{F1T40.16}
\]

The full Siegel weight is now the alternating form
\(F^*b_U\), with the sign convention (F1T40.12).
Its rank is even. We prove, for \(r\ge2k+2\), that
\[
 \mathcal R_k^{<2k}=\sum_{\substack{\ell\subset U\\\dim\ell=1}}C_\ell,
       \qquad E:=\dim\mathcal R_k^{<2k}<q^{2k+2r(k-1)}.
 \tag{F1T40.17}
\]
The left side sums the irreducible modules of maximal
Siegel rank below \(2k\). All lines of the symplectic
\(U\) are isotropic, so the displayed dimension bound
follows from (F1T40.16) and fewer than \(q^{2k}\) lines.

Here is the complete adaptation of the Fourier argument
to the actual connected orthogonal group. Take an
irreducible \(K\) of largest weight rank \(d<2k\)
and a nonzero full weight vector \(\Phi\) of that rank.
For a supported map \(F\), put
\[
 M=\operatorname{im}F,\quad R=M\cap M^\perp,
 \quad l=\dim M,\quad j=\dim R,
 \qquad d=l-j,\quad l+j\le2k.
 \tag{F1T40.18}
\]
The radical \(R\) is isotropic for the alternating form.
Take a two-dimensional subspace \(X_1\subset\ker F\)
and hold a complementary set of columns fixed.
The paired Fourier operator on \(X_1\) is an actual
\(H\)-operator: it is an \(SL(X)\)-conjugate of
(F1T40.14). One may choose the adapting basis in \(SL(X)\)
by adjusting its determinant on an unused complementary
line. The transformed slice has support only in
\(\operatorname{Hom}(X_1,R^\perp)\), since a vector
not orthogonal to \(R\) raises the polar rank of the
space containing \(M\) by at least two.
Fourier orthogonality therefore makes the original slice
invariant under \(\operatorname{Hom}(X_1,R)\).
Every kernel direction lies in such a two-plane, since
\(\dim\ker F\ge2\). Decomposing any translation into
these directions proves invariance under all
\(\operatorname{Hom}(\ker F,R)\). In composing these
translations, keep a fixed complementary subspace mapping
onto \(M\). Previously translated kernel coordinates lie
in \(R\), so the same slice-support argument remains
valid at each intermediate point.

In particular \(\Phi(F+\Delta)=\Phi(F)\) whenever
\(\Delta:X\to R\) and \(F|_{\ker\Delta}\) maps onto
\(M\): choose a complement to \(\ker F\) in
\(\ker\Delta\) and use the preceding actual paired
operators. For two full-image maps in the same branch,
the finite row-space calculation is
\[
 \Pr\bigl(F|_{\ker\Delta}\text{ not onto, or the
                  translated }F'\text{ map not onto}\bigr)
 \le2(q^l-1)(q^j-1)q^{-r}<1.
 \tag{F1T40.19}
\]
Indeed each prescribed nonzero row lies in the row space
of a uniform \(\Delta:X\to R\) with probability at
most \((q^j-1)q^{-r}\); the second translated map is
uniform as well. Here \(l+j\le2k\) and
\(r\ge2k+2\), so the bound is less than \(2q^{-2}\).
Thus all full-image coefficients on each branch agree,
including at \(q=2\).

Set \(j_0=(2k-d)/2\). A branch with \(j=j_0\) has
\(M=R^\perp\). For each such \(R\) and each
surjective quotient map \(\bar F:X\to U_R\) of the
given weight, let \(c_{R,\bar F}\) be the common
coefficient of its full-image lifts, and subtract exactly
\[
 c_{R,\bar F}q^{rj_0/2}i_R(e_{\bar F})
       =c_{R,\bar F}\sum_{\Delta:X\to R}e_{F+\Delta}.
 \tag{F1T40.20}
\]
Full-image lifts exist at this rank. The maximal branches
do not interfere, since inclusion of a full image
\(R'^\perp\) in \(R^\perp\), for equal dimensions of
the radicals, forces \(R=R'\).
The remainder \(\Phi'\) lies in the stable space
\(K+\sum_{\dim R=j_0}C_R\) of weight rank at most
\(d\), retains its full weight, and has no maximal branch.

If it is nonzero, choose a supported \(F\) of maximal
ordinary map rank and a two-plane \(X_1\subset\ker F\).
Its slice is supported on \(\operatorname{Hom}(X_1,R)\):
maximality puts added vectors in \(M\), and the weight
condition puts them in \(M^\perp\).
The established invariance makes this slice a nonzero
constant times that subspace's indicator. Its paired
Fourier transform is nonzero at every map
\(X_1\to R^\perp\). But \(R^\perp/R\) contains
\(M/R\) nondegenerately and an orthogonal complement
of dimension \(2k-2j-d\ge2\), since \(j<j_0\).
Map a basis of \(X_1\) to two vectors in that complement
with nonzero pairing. The transformed coefficient has
rank \(d+2\), a contradiction. Thus \(\Phi'=0\).
The generated irreducible lies in the code sum, and
\(C_R\subset C_\ell\) for any line \(\ell\subset R\).
The reverse inclusion in (F1T40.17) follows from the
quotient dimension in (F1T40.16). This proves the equality.

A minimal-level constituent of (F1T40.9) therefore has
weight rank \(2k\). Otherwise its projection from the
code sum is nonzero on some \(C_\ell\); the averaged
Hermitian adjoint embeds that constituent in the actual
\(\mathbf C[A^{k-1}]\), contrary to minimality.

\paragraph{4. Actual top projectors and all multiplicities.}
The full-rank fibre of a surjective \(F_0:X\to U\)
with fixed alternating pullback is an actual \(S\)-torsor.
Indeed the common radical is \(\ker F_0\), and every
other map is uniquely \(sF_0\), with \(s\in S\).
The stabilizer of its weight in the actual \(SL(X)\)
Levi surjects onto \(S\): act as any \(s\) on a
complementary quotient and identically on the kernel.
Its determinant is one because \(s\) is symplectic.
All full-rank weights are one \(SL(X)\)-orbit; an
arbitrary adapting \(GL(X)\)-matrix can have its
determinant corrected on the nonzero weight radical.

Write \(\mathcal R_k=\bigoplus_\tau E_\tau\otimes D_\tau\)
for the actual commuting \(S\)-action, and let
\(a=\dim E_\tau\), \(s=|S|<q^{4k^2}\).
The fixed-weight space of \(D_\tau\) is the irreducible
right-regular constituent \(E_\tau^*\) under the
stabilizer just exhibited. Every top-rank irreducible
\(H\)-constituent contains that weight. Hence
\(D_\tau\) has exactly one top constituent
\(\eta_\tau\), with multiplicity one; its other part
\(B_\tau\) has lower weight rank.
The actual evaluation map
\(\operatorname{Hom}_S(E_\tau,\mathcal R_k)\to\mathcal R_k\)
at a nonzero vector of \(E_\tau\) is injective, because
that vector generates the irreducible \(S\)-module.
Thus \(\dim B_\tau\le E\). Every selected original
minimal-level constituent is one of these \(\eta_\tau\).
The exact projector gives
\[
 \operatorname{Tr}_{D_\tau}(x)
      ={1\over s}\sum_{t\in S}\overline{\tau(t)}
                                  \operatorname{Tr}\mathcal R_k(x,t).
 \tag{F1T40.21}
\]
This is obtained by dividing the isotypic-projector trace
by \(a\), retaining its multiplicity space.

\paragraph{5. The uniform trace estimate on the literal commutator.}
Translation and multiplication operators on
\(\mathbf C[\operatorname{Hom}(X,U)]\) form an
orthogonal Weyl basis of its endomorphism space,
labelled by \(A\otimes U\). The explicit generators
(F1T40.12)--(F1T40.13), and the permutations by \(S\),
conjugate those labels by \(x\otimes t\), with unit
phases allowed. Equivalently this follows by transporting
the translation/multiplication basis through (F1T40.10);
the negative sign in (F1T40.12) gives exactly the natural
lower Siegel matrix on the labels. Therefore
\[
 |\operatorname{Tr}\mathcal R_k(x,t)|^2
                 \le q^{\dim\ker(x\otimes t-I)}.
 \tag{F1T40.22}
\]
For \(t\ne I\), its actual basis permutation also gives
\(\operatorname{Tr}\mathcal R_k(I,t)\le Dq^{-r}\).
For \(h\) in (F1T40.6)--(F1T40.7),
\[
 |\operatorname{Tr}\mathcal R_k(h,I)|\le Dq^{-2k},
 \qquad
 |\operatorname{Tr}\mathcal R_k(h,t)|\le Dq^{-r+2}
                               \quad(t\ne I).
 \tag{F1T40.23}
\]
The first inequality uses its natural rank at least two.
For the second, the restriction to
\(\ker(h-I)\otimes U\) has rank at least
\((2r-4)\operatorname{rank}(t-I)\ge2r-4\).

Set \(B=aD/s\). Subtracting the actual lower part
using (F1T40.17), and applying (F1T40.21), yields
\[
\begin{split}
 \dim\eta_\tau&>
 B\bigl(1-q^{4k^2-r}-q^{4k^2+2k-2r}\bigr),\\
 |\operatorname{Tr}\eta_\tau(h)|&<
 B\bigl(q^{-2k}+q^{4k^2-r+2}+q^{4k^2+2k-2r}\bigr).
\end{split}
 \tag{F1T40.24}
\]
For \(r\ge4k^2+9\), both the denominator error and
lower-part error are bounded by the corresponding T39
bounds, while the middle numerator term is at most
\(2^{-7}\). Thus
\[
 \frac{|\operatorname{Tr}\eta_\tau(h)|}{\dim\eta_\tau}
 <\frac{67585}{261119}<\frac7{27}.
 \tag{F1T40.25}
\]
This holds for every original constituent and every
restriction multiplicity. Summing gives
\[
             |\operatorname{Tr}_V\beta(h)|<7m/27.
\]
The original invariant Hermitian form and the literal
commutator (F1T40.6) give, exactly as in T39,
\[
 \operatorname{rank}(g-I)>5m/27\quad(g\notin Z).
 \tag{F1T40.26}
\]
Nonidentity scalars have rank \(m\). The original
\(b\ge12\) true inertia entries and exact budget
\(\sum_j\operatorname{rank}(g_j-I)=2m+6\)
then force \(m<27\), a contradiction. This proves
(F1T40.1), without assuming that a whole tensor factor
has the original target's cohomology.

\paragraph{6. Source, pairing and framework scope.}
The coefficient comparisons are the actual cover map,
root-element commutator, isotypic evaluation maps,
partial Fourier basis change, isotropic-code isometries,
paired operators in the root-generated group, and
averaged projector (F1T40.21). None changes the original
curve or its local system. All parabolic-cohomology,
primitive Jacobian-corner, Hermitian cup, restricted
polarization, Poincar\'e biextension, period and
Grothendieck--Adams comparisons therefore remain the
identity with pairing multiplier one. This is a strict
bound on original sources and adds no detector family.

The exact primary classification input is
Malle--Testerman, Theorem 24.24 and the retained natural
central-cover maps. The fixed-space parity and actual
level-generator character equality are
Guralnick--Larsen--Tiep, \emph{Inventiones Mathematicae}
235 (2024), Lemma 3.3, with its stated even-field
Guralnick--Tiep Lemma 5.8(ii) dependency.
The latter supplied proof cites the quasideterminant
identity in the \emph{Atlas}, p.~xii; that recursive
scope is recorded explicitly in the source audit.
All Fourier, code, top-projector and cutoff statements
used beyond those named inputs are proved above.
